% ============================================================
% Makros: Planspiel "Metro 2033: Die letzte Demokratie"
% Wird von beiden Hauptdateien (45 / 90 Minuten) eingebunden.
% Voraussetzung: \newif\ifkurz ist vorher definiert.
% ============================================================
\usepackage{graphicx}
\usepackage{array}
\usepackage{enumitem}
\usepackage[skins,breakable]{tcolorbox}
\usetikzlibrary{calc,shapes.geometric}

\emergencystretch=3em
\setlist[itemize]{leftmargin=1.3em,itemsep=1pt,topsep=2pt,parsep=0pt}
\setlist[enumerate]{leftmargin=1.6em,itemsep=1pt,topsep=2pt,parsep=0pt}

% ---------- Farben (je Gut / Gruppe) ----------
\definecolor{cN}{RGB}{46,125,50}   % Nahrung
\definecolor{cW}{RGB}{21,101,192}  % Wasser
\definecolor{cE}{RGB}{214,120,0}   % Energie
\definecolor{cM}{RGB}{198,40,40}   % Medikamente
\definecolor{cS}{RGB}{85,85,85}    % Schutzmaterial
\definecolor{warn}{RGB}{255,200,0}

% ---------- Ressourcensymbole ----------
\newcommand{\resdot}[2]{\tikz[baseline=(c.base)]\node[circle,fill=#1,text=white,inner sep=0pt,minimum size=3.8mm,font=\scriptsize\bfseries] (c) {#2};}
\newcommand{\rN}{\resdot{cN}{N}}
\newcommand{\rW}{\resdot{cW}{W}}
\newcommand{\rE}{\resdot{cE}{E}}
\newcommand{\rM}{\resdot{cM}{M}}
\newcommand{\rS}{\resdot{cS}{S}}
\newcommand{\phase}[1]{\tikz[baseline=(c.base)]\node[circle,draw,line width=0.6pt,inner sep=0pt,minimum size=4.2mm,font=\scriptsize\bfseries] (c) {#1};}

% ---------- Kästen ----------
\newtcolorbox{metrobox}[2][]{enhanced,breakable,colback=#2!6,colframe=#2!80!black,boxrule=0.9pt,arc=1.5mm,left=3mm,right=3mm,top=2mm,bottom=2mm,#1}
\newtcolorbox{lkbox}[1][]{enhanced,breakable,colback=black!4,colframe=black!65,boxrule=0.8pt,arc=1.5mm,left=3mm,right=3mm,top=2mm,bottom=2mm,#1}
\newtcolorbox{eventbox}[1][]{enhanced,breakable,colback=warn!14,colframe=black!85,boxrule=1.3pt,arc=1.5mm,left=3mm,right=3mm,top=2mm,bottom=2mm,#1}
\newtcolorbox{schneidbox}[2][]{enhanced,breakable,colback=white,frame hidden,boxrule=0pt,arc=0mm,left=3mm,right=3mm,top=2mm,bottom=2mm,coltitle=black,colbacktitle=#2!15,fonttitle=\bfseries,overlay={\draw[dashed,line width=0.9pt,#2!70!black](frame.south west)rectangle(frame.north east);},#1}

% Titelleiste einer Gruppenkarte: Farbe, Nummer, Name, Station
\newcommand{\gruppentitel}[5]{%
\begin{tcolorbox}[enhanced,colback=#1!85!black,colframe=#1!85!black,coltext=white,arc=1.5mm,boxrule=0pt,left=4mm,top=2.2mm,bottom=2.2mm]
{\Large\bfseries Gruppe #2 \,--\, #3}\hfill #5\\[0.8mm] Station: #4
\end{tcolorbox}\par\vspace{1mm}}

% Zahlentabelle einer Gruppe: je 5 Zellen (N & W & E & M & S)
\newcommand{\werte}[3]{%
\par\noindent
{\renewcommand{\arraystretch}{1.35}\setlength{\tabcolsep}{5pt}
\begin{tabular}{@{}l|c|c|c|c|c@{}}
 & \rN & \rW & \rE & \rM & \rS \\ \hline
Anfangsvorrat & #1 \\
Produktion pro Runde & #2 \\
Betriebskosten pro Runde & #3 \\
\end{tabular}}\par\vspace{1mm}}

% Marker (Ausschneidebogen): Farbe, Buchstabe, Beschriftung
\newcommand{\markerblock}[3]{%
\noindent{\bfseries #3}\par\vspace{1mm}\noindent
\foreach \i in {1,...,\markeranzahl}{%
 \tikz[baseline=(c.base)]\node[circle,draw=#1,fill=#1!12,line width=1.2pt,minimum size=10mm,text=#1!70!black,font=\normalsize\bfseries,inner sep=0pt] (c) {#2};\hspace{1.2mm}%
 \ifnum\i=\markerzeile\par\vspace{1.2mm}\noindent\fi}%
\par\vspace{3.5mm}}

% Bilanzblock für eine Runde (Gruppen-Buchführung)
\newcommand{\bilanzblock}[1]{%
\par\medskip\noindent{\bfseries #1}\par\smallskip
{\renewcommand{\arraystretch}{1}\setlength{\tabcolsep}{4pt}
\begin{tabular}{|p{6.4cm}|*{5}{>{\centering\arraybackslash}p{1.5cm}|}}
\hline
 & \rN & \rW & \rE & \rM & \rS \\ \hline
\rule{0pt}{6mm}Vorrat zu Rundenbeginn (inkl. Produktion) & & & & & \\ \hline
\rule{0pt}{6mm}Handel (+ erhalten, $-$ abgegeben) & & & & & \\ \hline
\rule{0pt}{6mm}Betriebskosten ($-$) & & & & & \\ \hline
\rule{0pt}{6mm}Krankheitskarte / Ereignis ($-$) & & & & & \\ \hline
\rule{0pt}{6mm}\textbf{Vorrat am Rundenende} & & & & & \\ \hline
\end{tabular}}\par}

% Stabilitätsleiste 0..10
\newcommand{\stabilitaetsleiste}{%
\begin{center}
\begin{tikzpicture}
\foreach \i in {0,...,10}{%
  \ifnum\i<1 \def\fc{red!55}\else\ifnum\i<5 \def\fc{orange!55}\else\ifnum\i<8 \def\fc{yellow!60}\else\def\fc{green!55!black!35}\fi\fi\fi
  \draw[fill=\fc,line width=1pt] ({\i*1.58},0) rectangle ++(1.58,1.7);
  \node[font=\Huge\bfseries] at ({\i*1.58+0.79},0.85) {\i};
}
\end{tikzpicture}
\end{center}}

% Schreiblinien, Checkbox
\newcommand{\linien}[1]{\par\vspace{1mm}\edef\linienzahl{#1}\letterbox[6.8mm]{\linienzahl}\par\vspace{2mm}}
\newcommand{\feld}{\Square\ }

% Marker-Anzahl je Variante
\ifkurz
  \newcommand{\markeranzahl}{18}\newcommand{\markerzeile}{9}
\else
  \newcommand{\markeranzahl}{24}\newcommand{\markerzeile}{12}
\fi

% ============================================================
% Eigene Vektorgrafiken (TikZ) -- frei von Fremdrechten
% Piktogramme: \iconN \iconW \iconE \iconM \iconS  {Farbe}{Skalierung}
% Banner:      \tunnelbanner     Raumplan: \raumplan
% ============================================================

% Pilz (Nahrung)
\newcommand{\iconN}[2]{\begin{tikzpicture}[scale=#2,baseline=(current bounding box.center)]
 \fill[#1] (-0.5,0.02) arc (180:0:0.5 and 0.46) -- cycle;
 \fill[#1,rounded corners=1.5pt] (-0.16,-0.46) rectangle (0.16,0.02);
\end{tikzpicture}}

% Tropfen (Wasser)
\newcommand{\iconW}[2]{\begin{tikzpicture}[scale=#2,baseline=(current bounding box.center)]
 \fill[#1] (0,0.5) .. controls (0.1,0.25) and (0.36,0.1) .. (0.36,-0.12) arc (0:-180:0.36) .. controls (-0.36,0.1) and (-0.1,0.25) .. (0,0.5);
\end{tikzpicture}}

% Blitz (Energie)
\newcommand{\iconE}[2]{\begin{tikzpicture}[scale=#2,baseline=(current bounding box.center)]
 \fill[#1] (0.12,0.5) -- (-0.3,-0.06) -- (-0.04,-0.06) -- (-0.14,-0.5) -- (0.32,0.1) -- (0.05,0.1) -- cycle;
\end{tikzpicture}}

% Kreuz (Medikamente)
\newcommand{\iconM}[2]{\begin{tikzpicture}[scale=#2,baseline=(current bounding box.center)]
 \fill[#1] (-0.16,-0.46) rectangle (0.16,0.46);
 \fill[#1] (-0.46,-0.16) rectangle (0.46,0.16);
\end{tikzpicture}}

% Schild (Schutzmaterial)
\newcommand{\iconS}[2]{\begin{tikzpicture}[scale=#2,baseline=(current bounding box.center)]
 \fill[#1] (-0.38,0.46) -- (0.38,0.46) -- (0.38,0.02) .. controls (0.38,-0.3) and (0.1,-0.42) .. (0,-0.5) .. controls (-0.1,-0.42) and (-0.38,-0.3) .. (-0.38,0.02) -- cycle;
\end{tikzpicture}}

% Tunnel-Banner (Breite 17,5 cm)
\newcommand{\tunnelbanner}{%
\begin{center}
\begin{tikzpicture}[x=1cm,y=1cm]
 \clip[rounded corners=2mm] (0,0) rectangle (17.5,3.8);
 \fill[black!92] (0,0) rectangle (17.5,3.8);
 \foreach \k in {0,...,6}{%
   \pgfmathsetmacro{\f}{1-0.13*\k}
   \pgfmathtruncatemacro{\pct}{88-9*\k}
   \fill[black!\pct!white,rounded corners=2mm] ({8.75-9.2*\f},{1.9-2.2*\f}) rectangle ({8.75+9.2*\f},{1.9+2.2*\f});
 }
 \draw[black!35!white,line width=1.4pt] (5.2,0) -- (8.55,1.75);
 \draw[black!35!white,line width=1.4pt] (12.3,0) -- (8.95,1.75);
 \foreach \y in {0.2,0.65,1.05,1.4}{%
   \pgfmathsetmacro{\xl}{5.2+(8.55-5.2)*\y/1.75}
   \pgfmathsetmacro{\xr}{12.3-(12.3-8.95)*\y/1.75}
   \draw[black!45!white,line width=1pt] ({\xl-0.35},\y) -- ({\xr+0.35},\y);
 }
 \foreach \f in {0.87,0.74,0.61,0.48}{%
   \fill[warn,opacity=0.3] ({8.75-9.2*\f*0.93},{1.9+2.2*\f*0.55}) circle ({0.3*\f});
   \fill[warn] ({8.75-9.2*\f*0.93},{1.9+2.2*\f*0.55}) circle ({0.11*\f});
   \fill[warn,opacity=0.3] ({8.75+9.2*\f*0.93},{1.9+2.2*\f*0.55}) circle ({0.3*\f});
   \fill[warn] ({8.75+9.2*\f*0.93},{1.9+2.2*\f*0.55}) circle ({0.11*\f});
 }
 \node[fill=warn,text=black,font=\bfseries\Large,inner xsep=4mm,inner ysep=1.8mm] at (8.75,3.05) {ENDSTATION 7};
 \node[text=white,font=\small] at (8.75,2.45) {Die letzte Demokratie};
\end{tikzpicture}
\end{center}}

% Raumplan (Klassenraum als Metrostation)
\newcommand{\raumplan}{%
\begin{center}
\begin{tikzpicture}[x=1cm,y=1cm,font=\footnotesize\bfseries]
 \draw[line width=1.2pt] (0,0) rectangle (17,7);
 \node[draw=cN!80!black,fill=cN!15,minimum width=4.6cm,minimum height=1.4cm,align=center] (g1) at (2.6,6.0) {Station 1\\Hydroponiker};
 \node[draw=cW!80!black,fill=cW!15,minimum width=4.6cm,minimum height=1.4cm,align=center] (g2) at (8.5,6.0) {Station 2\\Wasserwerker};
 \node[draw=cE!80!black,fill=cE!15,minimum width=4.6cm,minimum height=1.4cm,align=center] (g3) at (14.4,6.0) {Station 3\\Techniker};
 \node[draw=cM!80!black,fill=cM!15,minimum width=4.6cm,minimum height=1.4cm,align=center] (g4) at (2.6,1.0) {Station 4\\Sanitäter};
 \node[draw=cS!80!black,fill=cS!15,minimum width=4.6cm,minimum height=1.4cm,align=center] (g5) at (14.4,1.0) {Station 5\\Wächter (am Eingang)};
 \node[draw,ellipse,fill=warn!25,minimum width=5.6cm,minimum height=2.4cm,align=center] (konf) at (8.5,3.5) {Metro-Konferenz\\[-1pt]\footnotesize\mdseries Gemeinschaftsvorrat, Vertrag};
 \node[draw,fill=black!8,minimum width=3.6cm,minimum height=1.1cm,align=center] (bank) at (8.5,0.8) {Bank (Lehrkraft)};
 \foreach \g in {g1,g2,g3,g4,g5}{\draw[dashed,black!50] (\g) -- (konf);}
\end{tikzpicture}
\end{center}}

% ============================================================
% Atmosphäre: Zeilenabstand, Überschriften im Stil von Stationsschildern,
% Warnstreifen, Titelseite, eigene Motive (Tunnel, Strahlenzeichen, Gestalt mit Schutzmaske)
% Alle Motive sind eigene Zeichnungen (TikZ), keine Spielgrafiken.
% ============================================================
\usepackage{setspace}
\usepackage{etoolbox}

\definecolor{asphalt}{RGB}{28,28,30}
\definecolor{rust}{RGB}{150,62,25}
\definecolor{concrete}{RGB}{241,238,232}
\definecolor{glow}{RGB}{255,170,60}

\AtBeginDocument{\pagecolor{concrete}\setstretch{1.25}}
\AtBeginEnvironment{tabular}{\setstretch{1.0}}
\AtBeginEnvironment{longtable}{\setstretch{1.0}}
\addtokomafont{subsection}{\color{rust}}

% ---------- Abschnittsüberschrift (Stationsschild) ----------
\newcommand{\abschnitt}[1]{\par\vspace{1mm}%
\begin{tcolorbox}[enhanced,colback=asphalt,colframe=asphalt,coltext=white,boxrule=0pt,arc=0mm,left=8mm,top=2.5mm,bottom=2.5mm,borderline west={4mm}{0pt}{warn},borderline south={1.4pt}{0pt}{warn}]
{\Large\bfseries #1}
\end{tcolorbox}\par\vspace{1.5mm}}

% ---------- Titelleiste der Gruppenkarten (überschreibt makros.tex) ----------
\renewcommand{\gruppentitel}[5]{%
\begin{tcolorbox}[enhanced,colback=asphalt,colframe=asphalt,coltext=white,boxrule=0pt,arc=1mm,left=24mm,top=3mm,bottom=3mm,
 overlay={\fill[#1!85!black] (frame.south west) rectangle ([xshift=19mm]frame.north west);
          \node[white,font=\fontsize{30}{30}\selectfont\bfseries] at ([xshift=9.5mm]frame.west) {#2};}]
{\Large\bfseries #3}\hfill #5\\[0.8mm]\textcolor{warn}{Station: #4}
\end{tcolorbox}\par\vspace{1.5mm}}

% ---------- Warnstreifen ----------
\newcommand{\warnstreifen}[1][0.3]{\par\vspace{1mm}\noindent
\begin{tikzpicture}[x=1cm,y=1cm]
 \clip (0,0) rectangle (17.5,#1);
 \fill[warn] (0,0) rectangle (17.5,#1);
 \foreach \i in {-1,0,...,40}{\fill[asphalt] ({\i*0.8},0) -- ({\i*0.8+0.4},0) -- ({\i*0.8+0.4+#1},#1) -- ({\i*0.8+#1},#1) -- cycle;}
\end{tikzpicture}\par\vspace{1mm}}

% ---------- Strahlenzeichen (Dreiblatt) ----------
\newcommand{\trefcode}[1]{%
 \fill[#1] circle (0.13);
 \foreach \a in {90,210,330}{\fill[#1] ({\a-30}:0.22) -- ({\a-30}:0.85) arc (\a-30:\a+30:0.85) -- ({\a+30}:0.22) arc (\a+30:\a-30:0.22) -- cycle;}}
\newcommand{\trefoil}[2]{\begin{tikzpicture}[scale=#2,baseline=(current bounding box.center)]\trefcode{#1}\end{tikzpicture}}
\newcommand{\warnschild}[1]{\begin{tikzpicture}[scale=#1,baseline=(current bounding box.center)]
 \fill[warn] circle (1.05); \draw[asphalt,line width=1.6pt] circle (1.05);
 \begin{scope}[scale=0.95]\trefcode{asphalt}\end{scope}
\end{tikzpicture}}

% ---------- Gestalt mit Schutzmaske (Silhouette) ----------
\newcommand{\gestaltcode}[2]{%  Farbe, Linsenfarbe
 \fill[#1,rounded corners=5pt] (-0.75,-1.3) -- (-0.75,-0.5) -- (-0.3,-0.32) -- (0.3,-0.32) -- (0.75,-0.5) -- (0.75,-1.3) -- cycle;
 \fill[#1] (0,0.15) ellipse (0.5 and 0.62);
 \fill[#1] (-0.17,-0.62) rectangle (0.17,-0.3);
 \fill[#1] (0,-0.62) circle (0.2);
 \fill[#2] (-0.2,0.2) circle (0.14);
 \fill[#2] (0.2,0.2) circle (0.14);}
\newcommand{\gestalt}[3]{\begin{tikzpicture}[scale=#3,baseline=(current bounding box.center)]\gestaltcode{#1}{#2}\end{tikzpicture}}

% ---------- Tunnelszene (Breite 17,5 cm, Höhe frei) ----------
\newcommand{\szene}[1]{%
\begin{center}
\begin{tikzpicture}[x=1cm,y=1cm]
 \pgfmathsetmacro{\H}{#1}
 \pgfmathsetmacro{\cy}{0.56*\H}
 \clip[rounded corners=2mm] (0,0) rectangle (17.5,\H);
 \fill[black!95] (0,0) rectangle (17.5,\H);
 \foreach \k in {0,...,8}{%
   \pgfmathsetmacro{\f}{1-0.105*\k}
   \pgfmathtruncatemacro{\g}{5+8*\k}
   \fill[glow!\g!black,rounded corners=3mm] ({8.75-9.6*\f},{\cy-0.58*\H*\f}) rectangle ({8.75+9.6*\f},{\cy+0.58*\H*\f});
 }
 % Rohre an der Decke
 \foreach \s in {-1,1}{%
   \draw[black!75,line width=3pt] ({8.75-\s*8.75},{0.93*\H}) -- ({8.75-\s*0.6},{\cy+0.2});
   \draw[black!55,line width=1pt] ({8.75-\s*8.75},{0.90*\H}) -- ({8.75-\s*0.7},{\cy+0.12});
   \draw[black!70,line width=2pt] ({8.75-\s*8.75},{0.80*\H}) -- ({8.75-\s*0.9},{\cy+0.05});
 }
 % Gleise und Schwellen
 \pgfmathsetmacro{\vy}{\cy-0.25}
 \draw[black!30!white,line width=1.6pt] (4.3,0) -- (8.55,\vy);
 \draw[black!30!white,line width=1.6pt] (13.2,0) -- (8.95,\vy);
 \foreach \t in {0.06,0.17,0.3,0.46,0.66}{%
   \pgfmathsetmacro{\yy}{\t*\vy}
   \pgfmathsetmacro{\xl}{4.3+(8.55-4.3)*\t}
   \pgfmathsetmacro{\xr}{13.2-(13.2-8.95)*\t}
   \draw[black!40!white,line width=1.1pt] ({\xl-0.5},\yy) -- ({\xr+0.5},\yy);
 }
 % Gestalten im Gegenlicht
 \begin{scope}[shift={(8.2,{\cy-0.55})},scale=0.42]\gestaltcode{black}{glow!70!white}\end{scope}
 \begin{scope}[shift={(9.35,{\cy-0.5})},scale=0.36]\gestaltcode{black}{glow!70!white}\end{scope}
 % Wandlampen
 \foreach \f in {0.9,0.78,0.66,0.54,0.42}{%
   \foreach \s in {-1,1}{%
     \fill[warn,opacity=0.28] ({8.75+\s*9.6*\f*0.93},{\cy+0.58*\H*\f*0.45}) circle ({0.34*\f});
     \fill[warn] ({8.75+\s*9.6*\f*0.93},{\cy+0.58*\H*\f*0.45}) circle ({0.12*\f});
   }
 }
 % Warnschild an der Wand
 \begin{scope}[shift={(2.2,{\cy+0.25*\H})},scale=0.55]
   \fill[warn] circle (1.05); \draw[black,line width=1.6pt] circle (1.05);
   \begin{scope}[scale=0.95]\trefcode{black}\end{scope}
 \end{scope}
 % Nebel am Boden
 \shade[top color=black!0,bottom color=black!85] (0,{0.30*\H}) rectangle (17.5,0);
 % Stationsschild
 \node[fill=warn,text=black,font=\bfseries\Large,inner xsep=5mm,inner ysep=2mm] at (8.75,{\H-0.8}) {ENDSTATION 7};
\end{tikzpicture}
\end{center}}

% ---------- Titelseite ----------
\newcommand{\titelseite}{%
\szene{7.2}
\begin{center}
{\color{rust}\fontsize{34}{38}\selectfont\bfseries DIE LETZTE DEMOKRATIE\par}
\vspace{2mm}
{\large Politisches Planspiel in der Metrostation \textbf{Endstation 7}\par}
\end{center}
\warnstreifen[0.35]
\begin{center}
{\Large\bfseries \ifkurz Kurzvariante \,|\, 45 Minuten\else Vollvariante \,|\, 90 Minuten\fi}\\[1mm]
Spielleitung und Schülermaterial \,|\, inspiriert von der Idee der Inseldemokratie
\end{center}
\warnstreifen[0.35]
}

